\documentclass[12pt]{article}
\usepackage[T1]{fontenc}
\usepackage[utf8]{inputenc}
\usepackage{lmodern}
\usepackage{textcomp}
\usepackage{enumitem}
\usepackage{geometry}
\geometry{margin=1in}
\begin{document}
\begin{center}
Answer Key - Business Administration - Undergraduate Level Assessment 20 \
\small (Answers are ordered exactly as in the question paper.)
\end{center}

\section*{Multiple Choice}
\begin{enumerate}[label=\arabic*.]
\item \textbf{Answer:} B
\\ \textit{Options:} A) Price bundling.; B) Price elasticity.; C) Price inelasticity.; D) Price inflation.
\\ \textit{Note:} Price elasticity measures the responsiveness of the quantity demanded or supplied of a good or service to changes in its price, thereby illustrating the relationship between price and quantity.
\item \textbf{Answer:} A
\\ \textit{Options:} A) Expectancy; B) Instrumental; C) Classical; D) Contingency
\\ \textit{Note:} The Expectancy Theory posits that individuals make decisions based on their expectations of outcomes aligned with their values and beliefs, thereby guiding their behavior in a goal-oriented manner.
\item \textbf{Answer:} C
\\ \textit{Options:} A) Repeat purchases.; B) Buyphases.; C) Buyclasses.; D) Tenders.
\\ \textit{Note:} The term "buyclasses" refers to the three distinct categories of purchasing situations-new task, straight rebuy, and modified rebuy-that organizations encounter, as identified by Robinson, Faris, and Wind in their 1967 framework.
\item \textbf{Answer:} A
\\ \textit{Options:} A) Retaining the status quo; B) Understanding society; C) Harnessing diversity; D) Building capacity
\\ \textit{Note:} Retaining the status quo is not a characteristic of a corporate social responsibility framework because such frameworks aim to promote positive change and accountability in business practices rather than maintaining existing conditions that may be detrimental to society or the environment.
\item \textbf{Answer:} D
\\ \textit{Options:} A) International trade; B) Global economics; C) Global competition; D) Globalisation
\\ \textit{Note:} Globalisation refers to the interconnectedness of economies and markets across the world, where competition in one country is significantly affected by economic activities, policies, and competitive dynamics in other countries.
\end{enumerate}

\section*{True / False}
\begin{enumerate}[label=\arabic*.]
\setcounter{enumi}{5}
\item \textbf{Answer:} True
\\ \textit{Options:} A) True; B) False
\\ \textit{Note:} A succession plan is designed to ensure a smooth transition by identifying and preparing individuals to take on key roles when they become vacant, making it a formal process essential for organizational continuity.
\item \textbf{Answer:} True
\\ \textit{Options:} A) True; B) False
\\ \textit{Note:} The statement is true because the OECD highlights that when ownership is separated from control, managers may prioritize their own interests over those of shareholders, leading to corporate governance challenges.
\item \textbf{Answer:} False
\\ \textit{Options:} A) True; B) False
\\ \textit{Note:} Functional publics are not the primary source of authority and support for an organization's operations because authority typically derives from formal organizational structures and leadership rather than from specific functional groups or publics.
\item \textbf{Answer:} False
\\ \textit{Options:} A) True; B) False
\\ \textit{Note:} Qualitative research methods prioritize understanding human experiences and social phenomena through in-depth, open-ended questions and non-numerical data, rather than focusing on numerical data collection from many respondents.
\item \textbf{Answer:} True
\\ \textit{Options:} A) True; B) False
\\ \textit{Note:} A strategic business unit is designed to operate independently and manage its own strategy and resources, allowing it to concentrate on particular markets or products within the broader context of a corporation.
\end{enumerate}

\section*{Long Form Response}
\begin{enumerate}[label=\arabic*.]
\setcounter{enumi}{10}
\item \textbf{Answer:} Integrated Marketing Communications (IMC) is a strategic approach that seeks to unify all marketing communication tools, channels, and messages to ensure a consistent brand message across various platforms. By integrating diverse marketing efforts-such as advertising, public relations, social media, and direct marketing-IMC enhances brand recognition and loyalty. This consistency is crucial because it reinforces the brand's identity and values, making it easier for consumers to understand and connect with the brand. Moreover, IMC allows organizations to maximize their marketing impact by providing a cohesive narrative that resonates with target audiences, ultimately leading to increased engagement and sales. In today's fragmented media landscape, the role of IMC is more vital than ever in maintaining clarity and coherence in brand communication.
\\ \textit{Note:} correct because it effectively defines Integrated Marketing Communications (IMC) as a strategic approach that aligns various marketing channels to create a consistent brand message, emphasizing its importance for brand recognition and consumer understanding.
\item \textbf{Answer:} Broadcast media plays a crucial role in advertising by providing a dynamic platform for demonstrating product benefits and amplifying the effectiveness of marketing messages. Through visual and auditory elements, advertisements can showcase a product's features, usability, and emotional appeal, making it easier for consumers to understand its value. For instance, television commercials can depict products in real-life scenarios, allowing viewers to visualize their benefits in context. Additionally, the wide reach of broadcast media helps advertisers connect with diverse audiences, enhancing brand awareness and recall. Ultimately, the immersive nature of broadcast media not only captivates viewers but also reinforces the intended marketing message, leading to increased consumer engagement and, ultimately, sales.
\\ \textit{Note:} correct because it emphasizes how broadcast media utilizes visual and auditory elements to effectively illustrate product benefits and enhance marketing messages by placing them in relatable contexts for consumers.
\end{enumerate}

\vfill
\noindent\textit{End of Answer Key}
\end{document}